\section{Reviewed baseline, motivation, and scope}
\label{sec:baseline}

\subsection{The maintained ProveIt pipeline}

The reviewed source is the public ProveIt repository at commit
\hashcode{23893144e3aafd074a23cfb633222ab7e03e6d2e}.  At that commit, the project
combines cheap diagram reductions and obstructions with exact Khovanov/Bar--Natan
computation.  The commit message records 431 maintained tests and states
explicitly that general quasi-polynomial recognition remains unproved.  This
report accepts that boundary \cite{ProveIt2026}.

The relevant preprocessor is in
\code{Topology/UnknotRecognition/fast/fastunknot/simplify.py}, reviewed blob
\hashcode{2fcaabb51eaa45cf84bdee32cd1b4e18acf2c478}.  The file represents a diagram
by a fixed list of crossing darts and a mutable edge involution.  It performs
crossing-decreasing RI/RII moves incrementally.  When no such move remains, it
can trial an RIII inversion of a triangular face, keep a short sequence if the
sequence exposes RI/RII, and otherwise undo it exactly.

The existing recursive helper has three important properties.
\begin{enumerate}
\item It is \emph{sound}: successful RIII trials are retained only when a legal
      RI/RII move appears, and the resulting trace is independently replayable.
\item It is \emph{budgeted}: the total number of trial moves is capped per input
      crossing, because failed deep searches can dominate runtime.
\item It is \emph{path-local}: after the first RIII move, later candidates are
      drawn from a neighborhood of the immediately preceding move, rather than
      from the accumulated support of the entire trace.
\end{enumerate}

The synthesis reports document why this matters.  On a corpus of difficult
unknot survivors, depth two, three, and four searches simplified many more
instances than one-step lookahead, but a failed depth-four search on a
triangle-rich braid was much more expensive.  The current default therefore
uses shallow depth and a strict budget.  These data motivate a theorem that
separates global branching from genuinely local continuation.

\subsection{Relation to known Reidemeister bounds}

Hass and Lagarias proved an exponential upper bound on the number of
Reidemeister moves needed to trivialize an $n$-crossing unknot diagram
\cite{HassLagarias2001}.  Dynnikov proved monotonic simplification for
arc-presentations of the unknot \cite{Dynnikov2006}; Henrich and Kauffman used
that framework to bound diagram complexity and move counts
\cite{HenrichKauffman2011}.  Lackenby later proved that an $n$-crossing unknot
diagram can be reduced using at most $(236n)^{11}$ Reidemeister moves while the
crossing number remains at most $(7n)^2$ along the sequence
\cite{Lackenby2015}.

Those theorems establish polynomial-size certificates and powerful global
control.  They do not assert that a sequence can be partitioned into bounded
runs of RIII moves separated only by crossing-\emph{decreasing} RI/RII moves.
The present search deliberately forbids crossing increase.  Consequently, the
new parameterized theorem is neither a replacement for Lackenby's theorem nor a
consequence of it.

Lackenby's 2026 hierarchy algorithm gives a different route to unknot detection,
through incompressibility and essential hierarchies \cite{Lackenby2026}.  Its
abstract emphasizes avoidance of certain large searches, but the preprint does
not claim the quasi-polynomial running-time theorem sought by this project.  The
local RIII analysis here is orthogonal and can coexist with hierarchy-based
certificates.

\subsection{Relation to the exact fallback}

Khovanov homology detects the unknot \cite{KronheimerMrowka2011}, and Bar--Natan's
tangle/cobordism formalism supports local scanning computations
\cite{BarNatan2005}.  ProveIt exploits this exact route after cheaper stages
fail.  Report 25 changes none of that algebra.  Its role is to make one
preprocessing search asymptotically accountable and more complete under explicit
parameters.

\subsection{Methodological scope}

The OpenAI mathematics repository contains manuscripts and proof artifacts at
multiple verification levels, explicitly warning that unformalized outputs may
contain errors \cite{OpenAIMath2026}.  This report follows the useful part of
that model: theorem statements are separated from experiments, failed proof
ideas are recorded, and the integration patch remains opt-in.  It does not rely
on repository size or model provenance as mathematical evidence.

\begin{warning}[Novelty and verification]
The theorems below are new relative to the reviewed ProveIt source and its
synthesis, but no claim of priority in the wider mathematical literature is made
without expert review.  The proofs are human-readable and computationally
stress-tested, not formally verified.  The ProveIt adapter was not executed
inside a full checkout in this environment.
\end{warning}
